% LaTeX Template für Abgaben an der Universität Stuttgart
% Autor: Sandro Speth
% Bei Fragen: Sandro.Speth@iste.uni-stuttgart.de
%-----------------------------------------------------------
% Hauptmodul des Templates: Hier können andere Dateien eingebunden werden
% oder Inhalte direkt rein geschrieben werden.
% Kompiliere dieses Modul um eine PDF zu erzeugen.

% Dokumentenart. Ersetze 12pt, falls die Schriftgröße anzupassen ist.
\documentclass[12pt]{scrartcl}
\usepackage{booktabs}


% Einbinden der Pakete, des Headers und der Formatierung.
% Mit den \include und \input Befehlen können Dateien eingebunden werden:
% \include: Fügt einen Seitenumbruch nach dem Text ein
% \input: Fügt KEINEN Seitenumbruch nach dem Text ein
% LaTeX Template für Abgaben an der Universität Stuttgart
% Autor: Sandro Speth
% Bei Fragen: Sandro.Speth@iste.uni-stuttgart.de
%-----------------------------------------------------------
% Modul fuer verwendete Pakete.
% Neue Pakete einfach einfuegen mit dem \usepackage Befehl:
% \usepackage[options]{packagename}
\usepackage[utf8]{inputenc}
\usepackage[T1]{fontenc}
\usepackage[ngerman]{babel}
\usepackage{lmodern}
\usepackage{graphicx}
\usepackage[pdftex,hyperref,dvipsnames]{xcolor}
\usepackage{listings}
\usepackage[a4paper,lmargin={2cm},rmargin={2cm},tmargin={3.5cm},bmargin = {2.5cm},headheight = {4cm}]{geometry}
\usepackage{amsmath,amssymb,amstext,amsthm}
\usepackage[lined,algonl,boxed]{algorithm2e}
% alternative zu algorithm2e:
%\usepackage[]{algorithm} %counter mit chapter
%\usepackage{algpseudocode}
\usepackage{tikz}
\usepackage{hyperref}
\usepackage{url}
\usepackage[inline]{enumitem} % Ermöglicht ändern der enum Item Zahlen
\usepackage[headsepline]{scrlayer-scrpage} 
\pagestyle{scrheadings} 
\usetikzlibrary{automata,positioning}

% LaTeX Template für Abgaben an der Universität Stuttgart
% Autor: Sandro Speth
% Bei Fragen: Sandro.Speth@iste.uni-stuttgart.de
%-----------------------------------------------------------
% Modul beinhaltet Befehl fuer Aufgabennummerierung,
% sowie die Header Informationen.

% Überschreibt enumerate Befehl, sodass 1. Ebene Items mit
\renewcommand{\theenumi}{(\alph{enumi})}
% (a), (b), etc. nummeriert werden.
\renewcommand{\labelenumi}{\text{\theenumi}}


% Counter für das Blatt und die Aufgabennummer.
% Ersetze die Nummer des Übungsblattes und die Nummer der Aufgabe
% den Anforderungen entsprechend.
% Gesetz werden die counter in der hauptdatei, damit siese hier nicht jedes mal verändert werden muss
% Beachte:
% \setcounter{countername}{number}: Legt den Wert des Counters fest
% \stepcounter{countername}: Erhöht den Wert des Counters um 1.
\newcounter{sheetnr}
\setcounter{sheetnr}{42}

\newcounter{exnum}

% Befehl für die Aufgabentitel
\newcommand{\exercise}[1]{\section*{Excercise \theexnum\stepcounter{exnum}: #1}} % Befehl für Aufgabentitel

% Formatierung der Kopfzeile
% \ohead: Setzt rechten Teil der Kopfzeile mit
% Namen und Matrikelnummern aller Bearbeiter
\ohead{Muhammed Talha Sevinc (3723559)}
% \chead{} kann mittleren Kopfzeilen Teil sezten
% \ihead: Setzt linken Teil der Kopfzeile mit
% Modulnamen, Semester und Übungsblattnummer
\ihead{Bachelor Ringvorlesung\\
Winter Term 2023\\
Übungsblatt \thesheetnr}


% Counter für das Blatt und die Aufgabennummer.
% Ersetze die Nummer des Übungsblattes und die Nummer der Aufgabe
% den Anforderungen entsprechend.
% Definiert werden die Counter in FormatAndHeader.tex
% Beachte:
% \setcounter{countername}{number}: Legt den Wert des Counters fest
% \stepcounter{countername}: Erhöht den Wert des Counters um 1.
% \setcounter{sheetnr}{0} % Nummer des Übungsblattes
\setcounter{exnum}{1} % Nummer der Aufgabe

% Beginn des eigentlichen Dokuments
\begin{document}
% Nutze den \exercise{Aufgabenname} Befehl, um eine neue Aufgabe zu beginnen.
% Möchtest du eine Aufgabe in der Nummerierung überspringen, schreibe vor der Aufgabe: \stepcounter{exnum}
% Möchtest du die Nummer einer Aufgabe auf eine beliebige Zahl x setzen, schreibe vor der Aufgabe: \setcounter{exnum}{x}

\exercise{}
\begin{table}[!h]
    \centering
        \begin{tabular} {| c | c | c | c |}
            \hline Name & capital city & population density & official language(s)\\
            \hline Turkey & Ankara & 112 per $km^2$ & Turkish\\
            \hline Germany & Berlin & 239 per $km^2$ & German\\
            \hline France & Paris &  118 per $km^2$ & French\\
            \hline Denmark & Copenhagen & 138 per $km^2$ & Danish\\
            \hline Italy & Rome & 195 per $km^2$ & Italian\\
            \hline
        \end{tabular}
        \caption{Table of 5 countries with 4 specifications}
    \label{tab:CountryTable}
\end{table}

\exercise{}
    \begin{figure}[t]
        \centering
        \includegraphics[width=150pt]{wombat}
        \caption{Funny Wombat}
        \label{fig:Wombat}
    \end{figure}

\exercise{}
    \begin{figure}[!h]
        \centering
            \begin{tikzpicture}
                \node[state, initial] (A) {$A$};
                \node[state] (B) [below right=of A] {$B$};
                \node[state] (C) [right=of B] {$C$};
                \node[state, accepting] [above right= of C] (D) {$D$};
                \path[->] (A) edge node[above]{$a$} (D);
                \path[->] (A) edge node[above]{$b$} (B);
                \path[->, loop left] (B) edge node[left]{$b$} (B);
                \path[->] (B) edge node[above]{$a$} (C);
                \path[->,loop right] (C) edge node[right]{$a$} (C);
                \path[->] (C) edge node[above]{$b$} (D);

            \end{tikzpicture}
        \caption{Graphic for Problem 3}
        \label{fig:TikZ}
    \end{figure}

\exercise{}
    \cite{DBLP:journals/tifs/Frank0ASAK0RS23}
    \newline
    \cite{DBLP:journals/iacr/Groth16}

\exercise{}
    \begin{enumerate}
        \item Table of five Countries \ref{tab:CountryTable}\\
        \item Wombat Picture \ref{fig:Wombat}\\
        \item DFA \ref{fig:TikZ}
    \end{enumerate}


\cite{s2011}

\bibliographystyle{plain}
\bibliography{literature}


% Ende des Dokuments
\end{document}
